
% ============================================================================
% DOCUMENT 8: Physical AI Oncology Clinical Trial Site Premises Code
% ============================================================================

\part{Physical AI Oncology Clinical Trial Site Premises Code}
\setcounter{section}{0}

\section{Scope and Application}

These premises standards govern the physical site, grounds, security systems,
environmental controls, and operational conditions for Physical AI oncology
clinical trial facilities authorized under California law. The standards address
the unique premises requirements created by 24-hour autonomous robot operation,
patient-centric on-demand scheduling, and the deployment of 29 Physical AI
system instances across a 85,000 to 100,000 square-foot clinical
facility~\cite{sitespec2026}.

The premises code complements the building code standards (structural,
mechanical, electrical, technology) by addressing how the facility and its
grounds are maintained, secured, and operated during continuous clinical
operations. The 24-hour simulation demonstrated that sustained operations
require premises management that handles both minute-level details (individual
patient movements, robot state transitions, adverse event responses) and
facility-wide patterns across a full 24-hour cycle~\cite{newtrial2026}.

\section{Site Security}

\subsection{Physical Security Systems}

\begin{enumerate}[label=(\alph*)]
\item The premises shall be equipped with a comprehensive security system
  including the following components:
  \begin{enumerate}[label=(\arabic*)]
  \item Video surveillance covering all exterior entrances, parking areas,
    patient drop-off zones, lobby and reception areas, clinical wing corridors,
    and robot maintenance areas.
  \item Electronic access control on all exterior doors, all clinical wing
    entries, server rooms, pharmacy, robot maintenance bay, and any area
    housing controlled substances or sensitive equipment.
  \item Intrusion detection systems on all exterior doors, ground-floor windows,
    and roof access points.
  \item Duress alarm stations at the reception desk, each shift supervisor
    station, and the security monitoring room.
  \end{enumerate}

\item Security systems shall operate continuously during the 24-hour operating
  period without interruption.

\item Security camera recordings shall be retained for a minimum of 90 days
  and made available to the department and law enforcement upon request.
\end{enumerate}

\subsection{Access Control Zones}

The premises shall be divided into the following access control zones:

\begin{enumerate}[label=(\alph*)]
\item \textbf{Public Zone:} Lobby, reception area, patient waiting areas,
  restrooms, and the covered patient drop-off canopy. Accessible to patients,
  visitors, and the public during operating hours.

\item \textbf{Patient Zone:} Clinical wings, examination rooms, procedure
  suites, recovery areas, and patient circulation corridors. Accessible to
  enrolled patients with valid credentials and accompanying persons.

\item \textbf{Clinical Zone:} Areas where Physical AI systems perform clinical
  procedures, including surgical suites, RT vaults, needle-placement suites,
  and steerable needle suites. Restricted to authorized clinical personnel,
  site safety officers, and patients undergoing procedures.

\item \textbf{Technical Zone:} Server rooms, robot maintenance bay,
  decontamination area, and equipment storage. Restricted to authorized
  technical personnel and site safety officers.

\item \textbf{Administrative Zone:} Staff offices, break rooms, training rooms,
  and supply storage. Restricted to site personnel.
\end{enumerate}

\subsection{Cybersecurity of Physical Premises}

\begin{enumerate}[label=(\alph*)]
\item Physical AI system control networks shall be accessible only from within
  the Technical Zone, with no wireless access points connected to robot control
  systems.

\item The MCP-PAI server infrastructure shall be housed in the Technical Zone
  server room with physical access limited to authorized cybersecurity and
  system administration personnel~\cite{iche6r3adapt}.

\item Network infrastructure components (switches, routers, fiber distribution
  panels) shall be housed in locked enclosures within the Technical Zone.

\item Visitor and patient wireless networks shall be physically isolated from
  clinical and Physical AI system networks.
\end{enumerate}

\section{Patient Access and Flow}

\subsection{Arrival and Reception}

\begin{enumerate}[label=(\alph*)]
\item The main patient entrance shall be accessible 24 hours per day, 7 days
  per week, consistent with the on-demand scheduling model that demonstrated
  patient arrivals across all 24 hours of the simulation, including 15
  arrivals during the peak hour (hour 09) and continued arrivals during
  overnight hours~\cite{newtrial2026}.

\item A covered patient drop-off zone shall accommodate a minimum of three
  vehicles simultaneously, with a covered walkway extending at least 20 feet
  to the building entrance.

\item The reception area shall include automated check-in kiosks that verify
  patient identity, confirm scheduled procedures, and direct patients to the
  appropriate clinical wing.

\item The reception area shall include seating for a minimum of 20 patients and
  accompanying persons, reflecting the simulation's peak concurrent on-site
  population of 28 patients~\cite{newtrial2026}.
\end{enumerate}

\subsection{Internal Patient Circulation}

\begin{enumerate}[label=(\alph*)]
\item Primary patient circulation corridors shall be a minimum of 8 feet wide
  to accommodate hospital beds, wheelchairs, rehabilitation exoskeletons, and
  mobile Physical AI systems simultaneously.

\item Wayfinding shall use a combination of visual signage, floor markings,
  digital displays, and optional auditory guidance compatible with the 10
  Physical AI system categories that patients may
  encounter~\cite{patientinstructions2026}.

\item Patient flow through the facility shall follow the patterns validated in
  the simulation, with separate intake, procedure, and recovery pathways to
  minimize cross-traffic and wait times~\cite{newtrial2026}.

\item Recovery areas adjacent to each procedure zone shall accommodate patients
  during the post-procedure observation periods specified in the patient robot
  instructions (ranging from 2 minutes for imaging to 60 minutes for steerable
  needle and surgical procedures)~\cite{patientinstructions2026}.
\end{enumerate}

\section{Environmental Standards}

\subsection{Exterior Grounds}

\begin{enumerate}[label=(\alph*)]
\item The site shall maintain landscaping and hardscaping consistent with the
  character of the surrounding neighborhood and the San Francisco Planning
  Code requirements for the applicable zoning district.

\item Exterior lighting shall provide a minimum of 5 foot-candles at all
  pedestrian pathways, parking areas, and building entrances, ensuring patient
  safety during the 24-hour operating period.

\item A minimum of 10 percent of the total lot area shall be dedicated to
  landscaped open space, which may include patient garden areas accessible from
  the Public Zone.

\item Signage shall clearly identify the facility as a Physical AI oncology
  clinical trial site, with the authorization certificate number displayed at
  the main entrance.
\end{enumerate}

\subsection{Interior Environment}

\begin{enumerate}[label=(\alph*)]
\item Patient waiting areas and recovery spaces shall maintain interior
  temperatures between 68 and 74 degrees Fahrenheit with ambient noise levels
  not exceeding 45 dBA.

\item Lighting in patient areas shall be adjustable between 200 and 500 lux to
  accommodate different times of day and patient comfort needs during 24-hour
  operations.

\item Clinical procedure areas shall maintain lighting appropriate to the
  specific Physical AI system requirements, ranging from high-intensity surgical
  lighting in surgical suites to reduced ambient lighting in imaging bays and
  RT vaults.

\item Air quality in all patient-accessible areas shall be maintained at a
  minimum of MERV 13 filtration with continuous monitoring of CO2, temperature,
  and humidity.
\end{enumerate}

\section{Robot Operational Zones}

\subsection{Robot Movement Corridors}

\begin{enumerate}[label=(\alph*)]
\item Designated robot movement corridors shall connect clinical wings to the
  robot maintenance bay, charging stations, and decontamination area.

\item Robot corridors shall be separated from primary patient circulation
  pathways where possible, using time-of-day scheduling where physical
  separation is not feasible.

\item Floor surfaces in robot corridors shall be smooth, level (maximum slope
  1:48), and free of thresholds or transitions that could impede autonomous
  robot navigation.

\item Robot corridors shall be equipped with proximity sensors and visual
  warning indicators (flashing lights) to alert human occupants when robots
  are in transit.
\end{enumerate}

\subsection{Robot Maintenance Bay}

\begin{enumerate}[label=(\alph*)]
\item The robot maintenance bay shall accommodate a minimum of three Physical
  AI systems simultaneously for scheduled maintenance, calibration, or repair.

\item The bay shall include dedicated workstations with electrical, compressed
  air, and network connections appropriate for each Physical AI system
  category.

\item Spare parts storage shall be maintained for critical components of all
  10 Physical AI system types, with inventory levels sufficient to support
  the maintenance schedule demonstrated in the simulation (where SURG-01
  maintenance was conducted from 22:15 to 01:59 without disrupting clinical
  operations)~\cite{newtrial2026}.

\item The maintenance bay shall be acoustically isolated from patient areas to
  allow maintenance activities during nighttime hours without disturbing
  patients in adjacent recovery areas.
\end{enumerate}

\section{Emergency Access and Evacuation}

\begin{enumerate}[label=(\alph*)]
\item Emergency vehicle access shall be provided from at least two directions,
  with a minimum 26-foot-wide fire access lane along the full length of the
  building.

\item A dedicated ambulance bay shall accommodate a minimum of two ambulances
  simultaneously, separate from the patient drop-off zone.

\item Evacuation plans shall account for the specific needs of oncology patients
  who may be connected to treatment equipment, recovering from procedures, or
  using rehabilitation exoskeletons at the time of evacuation.

\item Physical AI systems shall execute automated safe-shutdown and position
  procedures upon receiving evacuation signals, including retracting surgical
  arms, lowering treatment couches, and moving to designated safe positions
  to clear evacuation routes~\cite{patientinstructions2026}.

\item Evacuation drills shall be conducted quarterly, including scenarios
  specific to Physical AI system emergencies such as robot control system
  failure, surgical suite power loss during a procedure, and cybersecurity
  incidents requiring physical network isolation.
\end{enumerate}

\section{Waste Management}

\begin{enumerate}[label=(\alph*)]
\item Medical waste generated by Physical AI-assisted oncology procedures shall
  be managed in accordance with the California Medical Waste Management Act
  (Health and Safety Code \S~117600 et seq.).

\item Radioactive waste from radiotherapy operations shall be managed in
  accordance with applicable NRC and California Radiologic Health Branch
  requirements.

\item Robot decontamination waste shall be classified and managed based on the
  type of contamination (biological, chemical, or radioactive).

\item Dedicated waste staging areas shall be located in the Technical Zone with
  separate containers for biohazardous, radioactive, pharmaceutical, and
  general waste.

\item Waste collection shall be scheduled to avoid conflicts with patient
  arrival and departure patterns, using the patient flow data from the
  simulation to identify optimal collection windows~\cite{newtrial2026}.
\end{enumerate}

